\section{动作误差何时与任务完成一致}
\noindent\textbf{本节结论与推导路径。}
当参考专家可靠、部署状态被示范分布充分覆盖、动作误差能够泛化，
且任务对动作偏差的敏感程度受控时，较小的动作风险能够支持更好的任务成功率下界，
但不能保证两个策略的实际成功率按平均动作误差排序。
本节先用有限时域性能差分恒等式，将失败概率差写成部署分布下的任务优势累积；
再通过动作敏感性、分布覆盖及泛化条件，将其上界转换为可泛化动作误差的函数；
随后用完整覆盖下的反例及状态相关敏感性，说明任务更依赖哪些状态上的预测准确性。
现实中的有限数据、有限覆盖和有损观察难以满足理想保证，
因此下一节研究语义及几何监督能否改善任务相关表示，
再下一节研究更新观察和历史能否补足执行中的信息缺失。
三者处理不同缺口，不把新增模块当作自动恢复所有理想条件的手段。

\subsection{任务目标与分析假设}
任务在固定时域 $T$ 内运行，成功事件为 $\Psi=1$，失败包含超时和预先指定的违规条件。
所有策略使用同一初始分布、环境转移和成功判据：
\begin{equation}
 J(\pi)=\Prob_\pi(\Psi=1),\qquad F(\pi)=1-J(\pi).
 \label{eq:success}
\end{equation}
为建立条件性结果，令 $x_t$ 为共同的充分决策状态，$\pi(x_t)$ 为实际执行指令。
必要的任务记忆、计划和时序状态可纳入 $x_t$，使终端事件可由增广状态表达。
这不声称单帧RGB满足Markov性。若只能使用有损观测或记忆编码，
其信息缺失须另行计入，不能在应用理论时省略。

参考策略 $\pi_E$ 假设为定义在所考察状态上的确定性专家。
参考策略不必全局最优，但须有可接受的 $J(\pi_E)$。
示范存在互相矛盾的动作模式时，动作条件均值不必仍是有效专家，
此时不能把平方误差最优解直接代入下面的专家假设。
令 $D$ 为固定的正对角动作尺度矩阵，定义
\begin{equation}
 e_\pi(x)=D^{-1}(\pi(x)-\pi_E(x)).
\end{equation}
$D$ 来自训练集，所有比较沿用相同尺度；归一化并不代表各关节具有相同任务代价。
以下讨论执行时刻的动作误差，不默认全动作块各槽位平均误差等于实际执行误差。

令 $Q_t^E(x,a)$ 表示在 $x$ 执行 $a$ 后始终遵循专家的终端失败概率，
$V_t^E(x)=Q_t^E(x,\pi_E(x))$，$A_t^E=Q_t^E-V_t^E$。
$d_\pi^t$ 为执行策略 $\pi$ 的状态分布，$\bar d_\pi$ 为 $(t,x)$ 的均匀时间混合。
由性能差分恒等式的有限时域形式，
\begin{equation}
 F(\pi)-F(\pi_E)
 =\sum_{t=0}^{T-1}\E_{d_\pi^t} A_t^E(x,\pi(x)).
 \label{eq:pd}
\end{equation}
附录给出证明。等式的对象是\emph{闭环访问分布下、按任务后果衡量的误差}，
不是训练集的平均动作距离。

\subsection{覆盖、泛化与任务敏感性的共同作用}
记专家分布平方动作风险为
$\epsilon_E=\E_{\bar d_E}\norm{e_\pi}^2$。
假设部署分布被专家分布控制，且动作偏离引起的任务风险有界：
\begin{align}
 \bar d_\pi&\ll\bar d_E,\quad
 \frac{d\bar d_\pi}{d\bar d_E}\le C, \label{eq:coverage}\\
 A_t^E(x,a)&\le L\norm{D^{-1}(a-\pi_E(x))}. \label{eq:sensitivity}
\end{align}
$C$ 表示覆盖程度而非样本条数，$L$ 表示任务对动作偏差的敏感程度。
覆盖假设要求相关状态及其频率，不只是物体初始位置网格齐全。
完整连续历史包含确定性动作时，两策略的历史分布甚至可能互相奇异；
因此不能未经验证就在全历史空间假定有限 $C$。
采用低维状态近似虽方便，也必须检查其充分性。
接触不连续、硬阈值成功事件或悬崖式失败附近，有限且有用的 $L$ 也可能不存在。

\begin{proposition}[动作风险到成功率的条件性界]
若式\eqref{eq:coverage}--\eqref{eq:sensitivity}成立，则
\begin{equation}
 J(\pi)\ge J(\pi_E)-TL\sqrt{C\epsilon_E}.
 \label{eq:successbound}
\end{equation}
若以至少 $1-\delta$ 的概率有
$\epsilon_E\le\widehat\epsilon_E+g_\delta$，则同概率下
\begin{equation}
 J(\pi)\ge J(\pi_E)-TL\sqrt{C(\widehat\epsilon_E+g_\delta)}.
 \label{eq:empbound}
\end{equation}
下界可截断到零。
\end{proposition}
由式\eqref{eq:pd}、Cauchy--Schwarz及密度比界立即得到，详见附录。
式\eqref{eq:empbound}中的 $g_\delta$ 不可忽略：
训练误差下降但泛化间隙扩大，不给出更强保证。
若仅有同轨迹相邻帧上的随机划分，不能把帧数当作独立样本数。
该界也可能很松，本文不以未知的 $C,L$ 生成数值成功率预测。

覆盖不全时，设
$\bar d_\pi=(1-\kappa)d_{\rm cov}+\kappa d_{\rm out}$，
其中 $d_{\rm cov}\le C\bar d_E$，式\eqref{eq:sensitivity}只在覆盖部分成立。
利用终端失败优势至多为1，
\begin{equation}
 F(\pi)-F(\pi_E)
 \le T\big[(1-\kappa)L\sqrt{C\epsilon_E}+\kappa\big].
 \label{eq:uncovered}
\end{equation}
覆盖域内误差趋零仍不能消除域外项。
这里 $\kappa$ 是未覆盖的平均访问质量，不是每条轨迹必然失败概率，
故保留保守的 $T$ 因子。更换策略也会改变 $C$ 和 $\kappa$。

\subsection{更小误差不等于更高成功率}
\emph{下界改善不等于两个真实成功率的单调排序。}
即使 $T=1$ 且训练测试分布完全一致，结论仍不成立。
设两个状态等概率出现，专家均输出零，状态1的失败概率为
$\min(1,100a^2)$，状态2任意动作均成功。
策略P分别输出 $(0.1,0)$，策略Q输出 $(0,0.2)$，则
\begin{equation}
 \epsilon_P=0.005<0.020=\epsilon_Q,\qquad
 J(P)=0.5<1=J(Q).
 \label{eq:counterexample}
\end{equation}
这不是覆盖缺失，而是误差发生位置及后果不同。
在实际任务中，工具接触前后的少量误差可能比自由运动误差重要得多。

严格的策略改进条件应相对于当前基线 $\pi_0$ 的任务优势描述：
\begin{equation}
 \sum_t\E_{d_{\pi_1}^t}A_t^{\pi_0}(x,\pi_1(x))<0
 \quad\Longrightarrow\quad J(\pi_1)>J(\pi_0).
 \label{eq:strict}
\end{equation}
普通L1/MSE没有直接监督该条件。
只有额外建立代理损失与任务优势的关系，并控制访问分布变化，才可谈一致性。
本文的合理主张是：\emph{降低可泛化的、任务相关的执行误差，是提高成功率的支持路径，
而非充分的模型排名规则。}

\subsection{从理想保证到实际缺口}
前面的保证指出三个现实困难。
\emph{有限数据与表征学习困难：}即使某类控制关系在示范中出现，
有限样本也未必足以使模型从高维外观中发现并保留该关系；低训练误差不等于低泛化误差。
\emph{覆盖不全：}外观组合缺失与控制状态缺失不同，
前者可能通过不变先验缓解，后者不能仅靠重新编码已知图像补全。
\emph{观察与执行不同步：}生成动作时的观察可能不充分或已经过时，
即使图像表示准确，也不能表达尚未观察到的后续变化。
这些缺口分别引出表示监督与增量反馈，但尚未解释为何应强调某些语义关系。

为连接关键状态与表示选择，将式\eqref{eq:sensitivity}的统一常数细化为
非负的状态相关敏感性 $w(t,x)$：
\begin{equation}
 A_t^E(x,a)\le w(t,x)\norm{D^{-1}(a-\pi_E(x))}.
 \label{eq:statesensitivity}
\end{equation}
定义 $\epsilon_{E,w}=\E_{\bar d_E}[w\norm{e_\pi}^2]$，
$\bar w_\pi=\E_{\bar d_\pi}w$，假设这些量有限。
由加权Cauchy--Schwarz及同一覆盖条件，
\begin{align}
 F(\pi)-F(\pi_E)
 &\le T\E_{\bar d_\pi}[w\norm{e_\pi}]\nonumber\\
 &\le T\sqrt{\bar w_\pi C\epsilon_{E,w}}.
 \label{eq:weightedtaskbound}
\end{align}
取 $w\equiv L$ 即回到式\eqref{eq:successbound}。
这说明高敏感状态上的误差对该保证影响更大，
而不只是这些状态应在叙述上被称为关键。
这里 $w$ 是分析中的任务敏感性，不是现有训练的权重，
更不是SAM2标签质量分数；关键事件标记也不自动等于已估计的 $w$。

由此得到下一节的明确问题：\emph{如果高敏感状态下需要区分的控制关系，
能够由可靠语义几何表示，辅助监督能否减少这些状态上的表示与动作误差？}
关系在视觉上显著并不足够，还必须对动作有用。
若关键差异来自视觉不能区分的隐变量，或来自完全未见的控制规律，
语义表示并不能解决它。式\eqref{eq:weightedtaskbound}提供需要改善的风险对象，
而不是直接证明语义或QToken已经改善了它。
